% !TEX program = xelatex
% ============================================================================
% One Physics Book — 小学与初中物理（一至九年级）
% Simplified Chinese edition of one_physics_book_1_primary_middle_school.tex.
%
% Build:  latexmk one_physics_book_1_primary_middle_school_zh.tex
% Output: build/one_physics_book_1_primary_middle_school_zh.pdf
% ============================================================================
\documentclass[11pt,a4paper,twoside,openany]{book}

\newcommand{\booklang}{zh}
\usepackage{styles/onephysics}
% Give TeX room around unbreakable Latin and math runs in Chinese lines.
\setlength{\emergencystretch}{3em}

% Brand the running header: rosette + ONE-COURSE.COM (linking to the site)
% above the chapter mark and the book title, on even pages; page number at the
% foot. See \ombrandheader in styles/onephysics.sty.
\ombrandheader

% Print, at the end of every exercise, a hyperlinked "Solution p. N" pointer
% to its solution in the end-of-book appendix.
\omsolutionlinks

% Version and date default to empty / today, but a release build supplies a
% generated version.tex that overrides them (see .github/workflows/release.yml).
\newcommand{\bookdate}{\today}
\newcommand{\bookversion}{}
\IfFileExists{version.tex}{\input{version.tex}}{}

% This volume's line on the shared cover (frontmatter/titlepage.tex).
\newcommand{\bookline}{第 1 册：小学与初中物理}

\title{One Physics Book}
\author{One Physics Book 贡献者}
\date{\bookdate}

\begin{document}

\frontmatter
% One Course cover (mockup 81): pale blue-gray header band with the rosette
% sitting on its lower edge, ONE-COURSE.COM wordmark, four-color quad lines.
% Drawn as a full-page overlay: "remember picture" needs two LaTeX passes
% (latexmk handles the reruns). Coordinates in cm from the page's south west
% corner; the page is A4, 21 x 29.7.
%
% Text strings come from styles/lang/<lang>.tex (\bookmaintitle, …) and from
% the book entry file (\bookline).
\begin{titlepage}
  \begin{tikzpicture}[remember picture, overlay]
    \begin{scope}[shift={(current page.south west)}]
      % header band and rosette astride its lower edge
      \fill[ocPaleBlue] (0,24.6) rectangle (21,29.7);
      \ocQuadLine{0}{21}{24.36}{0.24}
      \ocRosette{10.5}{24.6}{2.4}{58}
      % wordmark, linking out to the site (\omSiteURL, styles/onephysics.sty).
      % \href would recolor it (hyperref's urlcolor), so the ink is restated
      % inside the link.
      \node[text=ocInk, font=\sffamily\bfseries\LARGE] at (10.5,20.6)
        {\href{\omSiteURL}{\textcolor{ocInk}{\textls[210]{ONE-COURSE.COM}}}};
      % title block
      \node[text=ocDarkBlue, font=\fontsize{40}{44}\selectfont\bfseries]
        at (10.5,16.4) {\bookmaintitle};
      \node[text=ocInk, font=\Large\itshape] at (10.5,14.75)
        {\booktagline};
      \node[text=ocInk, font=\normalsize, align=center, text width=13cm]
        at (10.5,13.4)
        {\bookdescription};
      \ocQuadLine{6.6}{14.4}{12.0}{0.22}
      % this volume's name (defined by each book's entry file)
      \ifdefined\bookline
        \node[text=ocDarkBlue, font=\Large\bfseries] at (10.5,11.15)
          {\bookline};
      \fi
      % credits, version and date
      \node[text=ocInk, font=\normalsize] at (10.5,6.2)
        {\bookcontributors};
      \node[text=ocInk, font=\small] at (10.5,5.55)
        {Benjamin Virrion \;\textperiodcentered\; One Course};
      \node[text=ocInk, font=\small] at (10.5,4.7)
        {\ifdefempty{\bookversion}{}%
          {\bookversionlabel\ \bookversion\;\textperiodcentered\;}\bookdate};
      % bottom band
      \ocQuadLine{0}{21}{0}{1.0}
    \end{scope}
  \end{tikzpicture}
  \null
\end{titlepage}

% Reference page (verso of the title page): where the book lives and its licence.
\thispagestyle{empty}
\vspace*{\fill}
\begin{center}
\omcolophonsource\\[2pt]
\url{https://github.com/bvirrion/one-physics-book}

\bigskip
\omcolophonlicense\\[2pt]
\url{https://creativecommons.org/licenses/by-nc-sa/4.0/}\\[6pt]
\omcolophoncredit
\end{center}
\vspace*{\fill}
\clearpage

\chapter{前言}

本书的愿望很简单：用一部前后连贯的物理书，涵盖从幼儿园到本科毕业的全部内容，
为世界各地的读者而写。

本书力求简洁而严谨：课程由定义、例子、命题和定理构成，只要在相应阶段可以理解，
就给出证明，随后是精心挑选的练习。不加证明给出的结果都明确标为“承认”。
所有练习的完整解答集中在书末。

每个学年构成本书的一个部分，每个部分再分为若干章。
学年越早，面向的读者越年幼：插图更多，演算步骤更详细，练习也更平缓。

\section*{如何阅读本书}

各陈述在每章内编号，并以颜色区分：定义为蓝色，定理为红色，命题及其同类为橙色，
方法（处理常见任务的简短步骤）为绿色。练习按难度分级，从直接应用的（$\star$）
到较难的（$\star\star\star$）。

\section*{参与贡献}

本书的源代码存放在一个公开的 git 仓库中：
\begin{center}
\url{https://github.com/bvirrion/one-physics-book}
\end{center}
欢迎各种贡献，包括新的章节、勘误、更好的证明和更多的练习；
详见仓库根目录下的 \texttt{CONTRIBUTING.md}。

% 本书图片的版权说明（仅第 1 册）。
% 机器可读的授权记录位于 images/book1/CREDITS.md；增加或删除照片时，请保持两个文件同步。
\chapter*{图片来源}

本书中的插图均为原创。照片依其各自的自由许可使用，谨向作者致谢：

\begin{itemize}
  \item \emph{吸住回形针的磁铁}（第 6 章）：Donald
        Olszewski，Wikimedia Commons，CC~BY~4.0。
  \item \emph{沙漏}（第 8 章）：John Morgan，Wikimedia
        Commons，CC~BY~2.0。
  \item \emph{月球正面}（第 11 章）：NASA/GSFC/Arizona
        State University（Lunar Reconnaissance Orbiter），公有领域。
  \item \emph{带地照的蛾眉月}（第 11 章）：
        G.~Donatiello，Wikimedia Commons，CC0。
  \item \emph{白昼中的凸月}（第 11 章）：Cbaile19，
        Wikimedia Commons，CC0。
  \item \emph{溴：一种有颜色的气体}（第 14 章）：James
        St.~John，Wikimedia Commons，CC~BY~2.0。
  \item \emph{鼓和鼓槌}（第 19 章）：Vladimir Morozov，
        Wikimedia Commons，CC~BY~2.0。
  \item \emph{指南针}（第 25 章）：Evan-Amos，Wikimedia Commons，
        公有领域。
  \item \emph{太阳}（第 26 章）：NASA/SDO (AIA)，公有领域。
  \item \emph{八大行星}（第 26 章）：NASA——
        水星：NASA/JHU-APL/Carnegie（MESSENGER）；金星：
        NASA/JPL（Magellan 雷达拼接图）；地球：NASA，Apollo~17；
        火星：NASA/USGS（Viking）；木星：NASA/JPL/Space Science
        Institute（Cassini）；土星：NASA/JPL/Space Science
        Institute（Cassini）；天王星和海王星：NASA/JPL
        （Voyager~2）。均为公有领域。
\end{itemize}

每张照片的完整来源链接和许可证 URL 均列于本书代码库中的
\texttt{images/book1/CREDITS.md}。

\bigskip
除照片外，第 1--5、7、9--10、12--13、15--18、20--22、24、27--29、31--34、36--40、42、44--45、48、50、53--55、59--61、63、66 以及 69--70 章中的四十五幅图为 AI 生成的插图，由 OpenAI 图像生成功能根据本书编辑撰写的提示词生成，并在收录前已审核其物理准确性。每幅生成图像的完整提示词均记录于代码库中的
\texttt{images/book1/ai/PROMPTS.md}。
\clearpage

\tableofcontents

\mainmatter

% --------------------------------------------------------------------------
% Parts — one per year.
% --------------------------------------------------------------------------
% ============================================================================
% Grade 1 (age 6-7) — first year of primary school.
% Age-adapted physics-flavored science (the old French primaire had no
% physics course): observation, light, astronomy, magnets.
%
% One chapter per file. Solutions live in solutions/, one file per chapter,
% same numbering. Chapter bodies are placeholders until written.
% ============================================================================
\part{\omstr{part.grade1}}

\ominput{grade-1}{01-five-senses}
\ominput{grade-1}{02-hot-and-cold}
\ominput{grade-1}{03-light-and-shadows}
\ominput{grade-1}{04-day-and-night}
\ominput{grade-1}{05-floating-and-sinking}
\ominput{grade-1}{06-magnets}

% ============================================================================
% Grade 2 (age 7-8) — second year of primary school.
% Age-adapted physics-flavored science: measurement, states of water,
% air, the Moon, first forces.
%
% One chapter per file. Solutions live in solutions/, one file per chapter,
% same numbering. Chapter bodies are placeholders until written.
% ============================================================================
\part{\omstr{part.grade2}}

\ominput{grade-2}{01-measuring-length}
\ominput{grade-2}{02-measuring-time}
\ominput{grade-2}{03-ice-water-steam}
\ominput{grade-2}{04-air-around-us}
\ominput{grade-2}{05-moon-in-the-sky}
\ominput{grade-2}{06-pushes-and-pulls}
\ominput{grade-2}{07-balance-equilibrium}

% ============================================================================
% Grade 3 (age 8-9) — third year of primary school.
% Age-adapted physics-flavored science: states of matter, temperature,
% a first circuit, levers, the Sun, sound, mass.
%
% One chapter per file. Solutions live in solutions/, one file per chapter,
% same numbering. Chapter bodies are placeholders until written.
% ============================================================================
\part{\omstr{part.grade3}}

\ominput{grade-3}{01-solids-liquids-gases}
\ominput{grade-3}{02-temperature-thermometers}
\ominput{grade-3}{03-first-electric-circuit}
\ominput{grade-3}{04-levers-and-scales}
\ominput{grade-3}{05-sun-days-seasons}
\ominput{grade-3}{06-sound-around-us}
\ominput{grade-3}{07-measuring-mass}

% ============================================================================
% Grade 4 (age 9-10) — fourth year of primary school.
% Age-adapted physics-flavored science: changes of state, circuits,
% weight, magnetism, astronomy, sound, light.
%
% One chapter per file. Solutions live in solutions/, one file per chapter,
% same numbering. Chapter bodies are placeholders until written.
% ============================================================================
\part{\omstr{part.grade4}}

\ominput{grade-4}{01-changes-of-state}
\ominput{grade-4}{02-conductors-insulators}
\ominput{grade-4}{03-series-circuits}
\ominput{grade-4}{04-weight-and-mass}
\ominput{grade-4}{05-magnets-and-compass}
\ominput{grade-4}{06-solar-system}
\ominput{grade-4}{07-how-sound-travels}
\ominput{grade-4}{08-light-straight-lines}

% ============================================================================
% Grade 5 (age 10-11) — last year of primary school.
% Age-adapted physics-flavored science: energy, circuits, machines,
% mirrors, the night sky, heat, speed, pressure.
%
% One chapter per file. Solutions live in solutions/, one file per chapter,
% same numbering. Chapter bodies are placeholders until written.
% ============================================================================
\part{\omstr{part.grade5}}

\ominput{grade-5}{01-everyday-energy}
\ominput{grade-5}{02-series-parallel-circuits}
\ominput{grade-5}{03-simple-machines}
\ominput{grade-5}{04-mirrors-reflection}
\ominput{grade-5}{05-stars-night-sky}
\ominput{grade-5}{06-heat-and-insulation}
\ominput{grade-5}{07-speed-distance-time}
\ominput{grade-5}{08-air-pressure-weather}

% ============================================================================
% Grade 6 (age 11-12) — first year of lower secondary school.
% Bridging year toward the old French college physique-chimie program:
% measurement, matter, motion, circuits, light, astronomy, density.
%
% One chapter per file. Solutions live in solutions/, one file per chapter,
% same numbering. Chapter bodies are placeholders until written.
% ============================================================================
\part{\omstr{part.grade6}}

\ominput{grade-6}{01-measurement-in-science}
\ominput{grade-6}{02-mass-and-volume}
\ominput{grade-6}{03-water-states}
\ominput{grade-6}{04-describing-motion}
\ominput{grade-6}{05-circuits-and-safety}
\ominput{grade-6}{06-light-propagation}
\ominput{grade-6}{07-sun-earth-moon}
\ominput{grade-6}{08-density-floating}

% ============================================================================
% Grade 7 (age 12-13) — second year of lower secondary school.
% Physics parts of the old French 5e program: matter, circuits, light,
% plus motion and sound for coherence of the series.
%
% One chapter per file. Solutions live in solutions/, one file per chapter,
% same numbering. Chapter bodies are placeholders until written.
% ============================================================================
\part{\omstr{part.grade7}}

\ominput{grade-7}{01-states-of-matter}
\ominput{grade-7}{02-volume-mass-density}
\ominput{grade-7}{03-circuits-series-parallel}
\ominput{grade-7}{04-short-circuits-safety}
\ominput{grade-7}{05-light-sources-propagation}
\ominput{grade-7}{06-shadows-eclipses}
\ominput{grade-7}{07-moon-phases}
\ominput{grade-7}{08-motion-average-speed}
\ominput{grade-7}{09-sound-production}

% ============================================================================
% Grade 8 (age 13-14) — third year of lower secondary school.
% Physics parts of the old French 4e program: laws of current and
% voltage, resistance, colors, lenses, the speed of light, sound.
%
% One chapter per file. Solutions live in solutions/, one file per chapter,
% same numbering. Chapter bodies are placeholders until written.
% ============================================================================
\part{\omstr{part.grade8}}

\ominput{grade-8}{01-electric-current-effects}
\ominput{grade-8}{02-intensity-ammeter}
\ominput{grade-8}{03-voltage-voltmeter}
\ominput{grade-8}{04-circuit-laws}
\ominput{grade-8}{05-ohms-law}
\ominput{grade-8}{06-colors-spectra}
\ominput{grade-8}{07-lenses-images}
\ominput{grade-8}{08-speed-of-light}
\ominput{grade-8}{09-sound-pitch-loudness}

% ============================================================================
% Grade 9 (age 14-15) — last year of lower secondary school.
% Physics parts of the old French 3e program: gravitation, motion and
% energy, alternating current, electric power.
%
% One chapter per file. Solutions live in solutions/, one file per chapter,
% same numbering. Chapter bodies are placeholders until written.
% ============================================================================
\part{\omstr{part.grade9}}

\ominput{grade-9}{01-gravitation}
\ominput{grade-9}{02-weight-vs-mass}
\ominput{grade-9}{03-uniform-varied-motion}
\ominput{grade-9}{04-kinetic-energy-safety}
\ominput{grade-9}{05-alternating-current}
\ominput{grade-9}{06-electric-power-energy}
\ominput{grade-9}{07-alternator}
\ominput{grade-9}{08-energy-conversions}
\ominput{grade-9}{09-atoms-to-galaxies}


% --------------------------------------------------------------------------
% Solutions to all exercises
% --------------------------------------------------------------------------
\appendix
\part{\omsolutionspart}
% Solutions to the Grade 1 exercises — one file per chapter.
\chapter{\omstr{part.grade1}}

\ominputsol{grade-1}{01-five-senses}
\ominputsol{grade-1}{02-hot-and-cold}
\ominputsol{grade-1}{03-light-and-shadows}
\ominputsol{grade-1}{04-day-and-night}
\ominputsol{grade-1}{05-floating-and-sinking}
\ominputsol{grade-1}{06-magnets}

% Solutions to the Grade 2 exercises — one file per chapter.
\chapter{\omstr{part.grade2}}

\ominputsol{grade-2}{01-measuring-length}
\ominputsol{grade-2}{02-measuring-time}
\ominputsol{grade-2}{03-ice-water-steam}
\ominputsol{grade-2}{04-air-around-us}
\ominputsol{grade-2}{05-moon-in-the-sky}
\ominputsol{grade-2}{06-pushes-and-pulls}
\ominputsol{grade-2}{07-balance-equilibrium}

% Solutions to the Grade 3 exercises — one file per chapter.
\chapter{\omstr{part.grade3}}

\ominputsol{grade-3}{01-solids-liquids-gases}
\ominputsol{grade-3}{02-temperature-thermometers}
\ominputsol{grade-3}{03-first-electric-circuit}
\ominputsol{grade-3}{04-levers-and-scales}
\ominputsol{grade-3}{05-sun-days-seasons}
\ominputsol{grade-3}{06-sound-around-us}
\ominputsol{grade-3}{07-measuring-mass}

% Solutions to the Grade 4 exercises — one file per chapter.
\chapter{\omstr{part.grade4}}

\ominputsol{grade-4}{01-changes-of-state}
\ominputsol{grade-4}{02-conductors-insulators}
\ominputsol{grade-4}{03-series-circuits}
\ominputsol{grade-4}{04-weight-and-mass}
\ominputsol{grade-4}{05-magnets-and-compass}
\ominputsol{grade-4}{06-solar-system}
\ominputsol{grade-4}{07-how-sound-travels}
\ominputsol{grade-4}{08-light-straight-lines}

% Solutions to the Grade 5 exercises — one file per chapter.
\chapter{\omstr{part.grade5}}

\ominputsol{grade-5}{01-everyday-energy}
\ominputsol{grade-5}{02-series-parallel-circuits}
\ominputsol{grade-5}{03-simple-machines}
\ominputsol{grade-5}{04-mirrors-reflection}
\ominputsol{grade-5}{05-stars-night-sky}
\ominputsol{grade-5}{06-heat-and-insulation}
\ominputsol{grade-5}{07-speed-distance-time}
\ominputsol{grade-5}{08-air-pressure-weather}

% Solutions to the Grade 6 exercises — one file per chapter.
\chapter{\omstr{part.grade6}}

\ominputsol{grade-6}{01-measurement-in-science}
\ominputsol{grade-6}{02-mass-and-volume}
\ominputsol{grade-6}{03-water-states}
\ominputsol{grade-6}{04-describing-motion}
\ominputsol{grade-6}{05-circuits-and-safety}
\ominputsol{grade-6}{06-light-propagation}
\ominputsol{grade-6}{07-sun-earth-moon}
\ominputsol{grade-6}{08-density-floating}

% Solutions to the Grade 7 exercises — one file per chapter.
\chapter{\omstr{part.grade7}}

\ominputsol{grade-7}{01-states-of-matter}
\ominputsol{grade-7}{02-volume-mass-density}
\ominputsol{grade-7}{03-circuits-series-parallel}
\ominputsol{grade-7}{04-short-circuits-safety}
\ominputsol{grade-7}{05-light-sources-propagation}
\ominputsol{grade-7}{06-shadows-eclipses}
\ominputsol{grade-7}{07-moon-phases}
\ominputsol{grade-7}{08-motion-average-speed}
\ominputsol{grade-7}{09-sound-production}

% Solutions to the Grade 8 exercises — one file per chapter.
\chapter{\omstr{part.grade8}}

\ominputsol{grade-8}{01-electric-current-effects}
\ominputsol{grade-8}{02-intensity-ammeter}
\ominputsol{grade-8}{03-voltage-voltmeter}
\ominputsol{grade-8}{04-circuit-laws}
\ominputsol{grade-8}{05-ohms-law}
\ominputsol{grade-8}{06-colors-spectra}
\ominputsol{grade-8}{07-lenses-images}
\ominputsol{grade-8}{08-speed-of-light}
\ominputsol{grade-8}{09-sound-pitch-loudness}

% Solutions to the Grade 9 exercises — one file per chapter.
\chapter{\omstr{part.grade9}}

\ominputsol{grade-9}{01-gravitation}
\ominputsol{grade-9}{02-weight-vs-mass}
\ominputsol{grade-9}{03-uniform-varied-motion}
\ominputsol{grade-9}{04-kinetic-energy-safety}
\ominputsol{grade-9}{05-alternating-current}
\ominputsol{grade-9}{06-electric-power-energy}
\ominputsol{grade-9}{07-alternator}
\ominputsol{grade-9}{08-energy-conversions}
\ominputsol{grade-9}{09-atoms-to-galaxies}


\backmatter
\printindex

\end{document}
